\chapter{Chess: The Closed Hidden-State Case}
\label{ch:chess}

\begin{quotation}
\noindent\textit{Synopsis.} Chess serves as the canonical closed
hidden-state domain. This chapter examines the gap between the visible
board and the additional state required for correct continuation, such as
castling rights and the repetition count, and shows that because this
hidden state is finite and enumerable, augmentation restores sufficiency in
principle. Chess offers a clean illustration of the framework's core
concepts without the complications introduced by open-ended real-world
domains, and the chapter closes by revisiting, through the case of a
physical board, the question of what direct access actually means.
\end{quotation}

\section{The Board as Substrate}

A chess position's visible state -- piece placement -- omits castling
rights, en passant eligibility, and the repetition count, each of which
is required to determine legal continuations in some positions. This is
the canonical witness for \cref{def:closed-open}'s closed case: the
omitted variables are drawn from a short, enumerable list, in sharp
contrast to painting's open causal space (Chapter~\ref{ch:painting}).

\begin{example}
\label{ex:castling-position}
Two positions can have identical piece placement -- the same pieces on
the same squares -- while differing in whether either player retains the
right to castle, if one position arose from a king or rook having moved
and moved back while the other did not. A rendering that records only
piece placement cannot distinguish these; a rendering that also records
castling rights can. This is a minimal pair in the exact sense of
\cref{def:minimal-pair}: two source states differing in one named
distinction, told apart in the source, and identified by a rendering
that omits it.
\end{example}

\section{Augmentation Worked in Full}

Augmenting the visible board with $H = \{\text{castling rights, en
passant status, repetition count}\}$ satisfies \cref{prop:closed-sufficiency}:
the augmented rendering $r' = (r, H)$ is sufficient for determining every
legal continuation, since $H$ exhausts the game's hidden state relevant
to legality. This can be stated even more precisely using
Appendix~\ref{app:proof}'s recovery-task machinery: each element of $H$
corresponds to a recovery task $\tau_h$ (\cref{def:recovery-task}) --
``does this side retain castling rights,'' say -- and
\cref{lem:recovery-sufficiency} gives directly that $r'$ is sufficient
for $\tau_h$ if and only if $h \notin \discard{r'}$. Since $H$ is
included in full, every such $\tau_h$ is satisfied simultaneously, which
is exactly what \cref{prop:closed-sufficiency}'s conclusion amounts to
once stated at the level of individual recovery tasks rather than only
at the level of the augmented rendering as a whole. Chess is, in this
sense, the domain where the recovery-task family $T_{\mathrm{rec}}$
introduced to construct Appendix~\ref{app:proof}'s witness happens to be
finite and known in advance, rather than open -- the one condition
under which \cref{prop:task-unavailability}'s impossibility result does
not bind.

\section{The Simul Player as a Degenerate Relay}

A simultaneous-exhibition player moving between boards is formalized as a
degenerate one-stage relay in which $B_0$ is deliberately reapplied fresh
at each encounter with a given board, rather than carried forward across
encounters: the player reads the current position cold and acts on it,
with no continuity of an internal working state between visits.

\begin{proposition}[Stateless Re-Entry Is Sufficient When the Board Is]
\label{prop:stateless-reentry}
If the augmented rendering $r' = (r, H)$ of a chess position is
sufficient for the task of determining the best legal move
(\cref{prop:closed-sufficiency}), then a player re-entering the game
statelessly at that position -- with no memory of how it arose -- can
determine the same best move as a player with full continuous memory of
the game.
\end{proposition}

\begin{proofsketch}
By \cref{prop:closed-sufficiency}, $r'$ is sufficient for the task, so
$\admis{\tau}{r'} \subseteq \mathcal{C}_\tau(e)$ regardless of which
consumer is evaluating it or what that consumer's history has been. A
continuously-present player's continuation is computed from
$\admis{\tau}{r'}$ exactly as a freshly re-entering player's is, since
neither player has access to information beyond $r'$ once castling
rights, en passant status, and repetition count are included in it. The
game's history matters only insofar as it determines $H$'s current
values, and $H$ is, by hypothesis, already reflected in $r'$.
\end{proofsketch}

\cref{prop:stateless-reentry} is what licenses treating the simul player
as a legitimate limiting case rather than a curiosity: the entire reason
walking a circle of boards and reading each one cold works at all is that
chess's hidden state is closed and, when augmented, fully captured in the
position itself, exactly as \cref{sec:coord-augmentation}'s discussion of
augmentation as a coordination mechanism requires. This contrasts sharply
with the genuine multi-stage relays of Chapter~\ref{ch:pipelines}, where
state \emph{is} intended to carry forward across stages that do not share
a single fixed task, and where \cref{thm:composition-failure}'s failure
mode becomes possible precisely because that carried-forward state is
uncoordinated across stages with differing priorities -- a condition that
does not arise in the simul case because every visit to a given board
shares exactly one task.

\section{Revisiting the Hard Test Case}

The physical board in front of a player is, per
\cref{oq:unmediated-access}, the strongest available candidate for
fidelity $\fidelity{r}{\tau} = 1$ under \cref{def:fidelity}. In the
vocabulary of Chapter~\ref{ch:mode-relativity}, the question is precisely
what $r_{c,t}$ is for a player $c$ at the moment of moving: is it the
board itself, or is it $c$'s visual percept of the board, already
processed by whatever mechanism governs perception? If the narrow reading
of \cref{ch:mem8} is correct, this distinction can be drawn cleanly and
$r_{c,t}$ can be treated as the board itself for practical purposes,
recovering the naive intuition that a player has direct access to the
game. If the broad reading is correct, $r_{c,t}$ is already a
reconstructive product of the player's perceptual system, and the
question of whether fidelity ever actually reaches $1$ here --- rather
than merely approaching it closely enough for chess's purposes --- is
left open pending resolution of \cref{oq:mem8-fork}. Nothing in this
chapter's treatment of augmentation and stateless re-entry depends on
settling that question: \cref{prop:stateless-reentry} holds given
whatever $r_{c,t}$ turns out to be, provided it includes $H$, and the
fork bears only on how literally ``the board itself'' can be taken as
that substrate.
